\section{从现实问题到概率模型}
\label{sec:05-Probability/01-Probability-Spaces/03}

你拿到一道概率题：``从十个人中随机选两人，两人都是女性的概率是多少？''已知十人中有四位女性。

你翻了翻前面的公理，公理没意见——它们只保证概率空间内部没有矛盾。但公理也没帮你回答一个更前置的问题：``随机选两人''到底对应什么样的概率空间？是十个人写在小纸条上，闭眼一次抓两张？还是先抓一张、不放回再抓一张？还是每个人独立地以某个概率被选中，结果可能选中零个人、一个人、三个人？

你可能会想：这些有什么区别？最后不都是两个人吗？我们来算算看。

\subsection{同一个题目，不同翻译，不同答案}

先试试不加思索的翻译。既然``选两人''，样本空间就是所有二人组合。总共有 \(\binom{10}{2}=45\) 种组合。其中两人都是女性的组合有 \(\binom{4}{2}=6\) 种。概率 \(6/45=2/15\approx0.133\)。

好，接下来换一种翻译。考虑顺序：先选一人，再选一人。样本空间变成有序对 \((a,b)\)，其中 \(a\neq b\)，共有 \(10\times9=90\) 种。其中两人都是女性的有序对有 \(4\times3=12\) 种。概率 \(12/90=2/15\)。答案一样？

这次一样是因为两种翻译下的``有利结果''和``全部结果''同时乘了同一个因子 \(2\)（有序 vs 无序）。但下面的翻译就不一样了。

把样本空间压缩成``选中女性的人数''：\(\{0,1,2\}\)。三个结果，其中选到两位女性只是三分之一？当然不对——三个结果不等可能。\(0\) 位女性的概率是 \(\binom{6}{2}/\binom{10}{2}=15/45=1/3\)，\(1\) 位是 \(\binom{4}{1}\binom{6}{1}/\binom{10}{2}=24/45=8/15\)，\(2\) 位是 \(6/45=2/15\)。它们不是各 \(1/3\)。

但凭什么默认``所有二人组合等可能''？题面只写了``随机选两人''，没写``等可能地选''。你的默认来自日常经验——大部分抽签、抽奖、随机分组确实设计成等可能的，因为那是最自然的``公平''——但这个默认是一个实质性的建模假设，不是逻辑必然。

现在把场景换一下。你在做一个推荐系统，要从十篇文章中选两篇推荐给用户，但你心里清楚：用户对科技类文章更有兴趣，科技类文章在你的模型里应该有更高的被推荐概率。此时``等可能''就不再是合适的假设。``随机''两个字没有规定概率质量怎么分配。

\subsection{两个机制，两个分布：放回 vs 不放回}

把刚才的思考推进一步。十个零件中有四个次品，你随机抽取零件检验。如果你每次抽完放回、重新混匀，抽取方式和不放回会产生完全不同的概率。用具体数字感受这个差距。

\textbf{放回抽样。}每次抽到次品的概率恒为 \(4/10=0.4\)。抽两次，两次都是次品的概率：\((0.4)^2=0.16\)。恰好一次次品的概率：\(\binom{2}{1}\times0.4\times0.6=0.48\)。零次次品的概率：\((0.6)^2=0.36\)。

\textbf{不放回抽样。}第一次次品概率 \(4/10=0.4\)。在第一次次品的条件下，第二次次品概率 \(3/9\approx0.333\)。所以两次都是次品：\(0.4\times(3/9)=0.133\)。两次都不是次品：\((6/10)\times(5/9)=0.333\)。恰好一次次品：剩下的 \(1-0.133-0.333=0.533\)。

对比两组数字：
\[
\begin{aligned}
\text{放回：}&P(\text{两次次品})=0.16,\quad P(\text{一次})=0.48,\quad P(\text{零次})=0.36.\\
\text{不放回：}&P(\text{两次次品})=0.133,\quad P(\text{一次})=0.533,\quad P(\text{零次})=0.333.
\end{aligned}
\]

不放回时，抽到两次次品的概率更低（因为抽走一个次品后，残局更干净了），抽到恰好一次次品的概率更高（同样的原因：抽走一个次品后，残局更干净，第二次更容易抽到好零件；反之，先抽到好零件，残局中次品比例升高，第二次更容易抽到次品——两种顺序恰好互补，推高了``恰好一次''的概率）。

如果零件总数是一万个而不是十个，放回和不放回的差距会缩小——拿走一个对整体比例的影响微乎其微。但零件只有十个时，差距大到不容忽视。你的概率模型必须明确选择其中一个机制，否则 \(0.16\) 和 \(0.133\) 都是``正确''的答案——只不过对应不同的实验设计。

\subsection{先问机制，再选公式}

``从一副牌中抽五张''至少可以意味着：

\begin{itemize}
\tightlist
\item 洗匀后一次性发五张——每手五张牌等可能；
\item 洗匀后每次抽一张，抽后放回洗牌——同一张牌可能重复出现；
\item 洗匀后每次抽一张，不放回——和你一次性拿五张的概率分配相同，但时间顺序被记录了；
\item 牌有磨损，某些牌更容易被抽到——古典概型的等可能假设崩塌。
\end{itemize}

第一种和第三种在概率分配上等价（如果你不关心顺序的话），但和第二、第四种完全不是同一个概率空间。题目里如果只写``随机抽五张牌''而不写有没有放回、牌是否均匀，你脑子里必须有报警声：这句话不完整。

建立概率模型的时候，可以沿着这五个问题逐项写下来：

\begin{enumerate}
\tightlist
\item \textbf{一次完整实验的结果是什么？}是一张牌？一手牌？还是一个带顺序的序列？把结果单元定清楚，样本空间才有边界。
\item \textbf{哪些子集被当作事件？}有限空间里通常全取，但如果你只关心某些统计量（比如点数之和），可以只考虑由这些统计量生成的事件族。
\item \textbf{概率怎么分配？}等可能？按权重？还是由一个更底层的物理过程（比如洗牌算法）导出？这里每做一个选择，都是在对现实做一个承诺。
\item \textbf{分配依赖了哪些现实假设？}牌是匀质的？洗牌充分随机？抽牌者没有偏好？把这些假设写下来，以后如果答案翻了，你知道去哪里找原因。
\item \textbf{结论对假设有多敏感？}如果把等可能换成轻微不均匀，概率变化是几个百分点还是直接翻倍？敏感的地方就是你该花力气验证的地方。
\end{enumerate}

每一步都不需要高深的数学。但跳过任何一步，后面的计算都可能变成``逻辑上正确但回答的是另一个问题''。

\subsection{计数公式是工具，不是概率模型}

你看到一道题，条件反射写出

\[
\frac{\binom{4}{2}\binom{48}{3}}{\binom{52}{5}},
\]

然后说``这就是从一副牌中抽五张，恰好两张 A 的概率''。数字算出来是 \(4\times3/2\times48\times47\times46/6\) 除以 \(52\times51\times50\times49\times48/120\)……约等于 \(0.04\)。没错吗？

分子和分母之间的除法成为概率，依赖于一个你写组合数时就默认了的前提：\textbf{所有五张牌的无序组合是等可能的}。这个前提藏在 \(\binom{52}{5}\) 这个分母里——它只是在数有多少种组合，不是概率本身。把计数变成概率，是你额外加入的``等可能''假设在起作用。

如果抽牌机制是有放回的，或者牌与牌之间有黏连（旧牌比新牌更容易粘在一起被同时抽出），组合数公式的分子和分母在计数上可能仍然正确，但它们相除之后不再等于概率——因为那些组合不再是等可能的。

计数工具在\hyperref[ch:027]{``组合计数：先描述对象怎样生成，再决定怎样数''}中有系统展开。这里只需要一个认知锚点：\textbf{计数给你分子和分母，概率模型决定分子除以分母是否有意义}。

\subsection{连续均匀：长度、面积、体积，不是``随便''}

``在 \([0,10]\) 上均匀随机选一个点，它落在 \([2,5]\) 内的概率是多少？''

你写出

\[
P(2\le X\le5)=\frac{5-2}{10}=0.3.
\]

直觉上没有问题——区间长度之比。但为什么是长度之比？因为``均匀''在这里的意思是：任意两个长度相等的子区间有相同的概率。从这个条件出发，可以推出概率必然与区间长度成正比。然后规范化公理把比例常数定为 \(1/10\)。

你可能会想：能不能定义一种``均匀''，让概率和长度的平方成正比？可以，但那就不再是你日常语言里说的``均匀''了——你定义的东西只是一个合法的概率测度，恰好具有某个函数形式。

平面区域上的均匀随机点按面积算，立体中的按体积算。每一次``均匀''都必须相对于某个测度（长度、面积、体积、计数测度……）来理解。没有抽象的均匀，只有针对特定结构的均匀。

Bertrand 弦悖论是这里最著名的警告。问题是这样：``在圆中随机取一条弦，弦长超过圆内接等边三角形边长的概率是多少？''

你可能会想：弦由两个端点确定，在圆周上均匀随机选两个点，计算弦长。这样算出来的概率是 \(1/3\)。

但换一种``随机取弦''的方式：在圆内均匀随机选一点作为弦的中点，然后做以该点为中点的弦。这样算出来的概率是 \(1/4\)。

再换一种：均匀随机选一个方向，再均匀随机选圆心到弦的距离，这样变成了 \(1/2\)。

三个答案都是对的——在各自的概率模型下。矛盾不在概率论自身，而在``随机取弦''这四个字没有把取弦的机制说完整。日常语言中的``随机''是一个高度压缩的词汇，它在不同的物理操作下展开成不同的概率空间。数学只能保证每个展开之后的推理不出错；展开之前的翻译工作，你得自己做。

\subsection{建模清单：拿到题目先别动手算}

前面几个陷阱指向同一个教训：概率计算的真正难点通常不在计算本身，而在动手之前的那几步翻译。下面是一份可以贴在你脑子里的自问清单。每拿到一个问题，逐条过一遍。

\textbf{第 1 步：结果单元。}一次完整实验产生的一整条记录长什么样？是掷一次骰子看点数，还是连掷十次看序列，还是掷到六就停看等待次数？把单个结果的结构定清楚，样本空间的边界才有意义。如果你发现自己在计算过程中偷偷改变了结果的定义——比如一开始说``看点数''，算到中间又按``看奇偶''来分类——那就是建模边界不稳的信号。

\textbf{第 2 步：等可能假设。}你打算让哪些结果等可能？这个决定直接锁定了概率分配的方式。如果题目写了``公平骰子''，你在借用物理世界的对称性。如果题目写了``随机抽取''，你假设了抽取机制是均匀的。每做一个等可能假设，都在心里用一句话写下你的理由；如果写不出理由，你就该怀疑这个假设是不是只在帮自己简化计算。

\textbf{第 3 步：放回还是不放回。}重复抽取时，每次抽取之后样本空间是否复原？用具体的残局数字检验——袋中有 \(5\) 球，第二次抽取的概率在放回下仍是原来的比例，在不放回下残局已变——而不是凭直觉说``差不多''。在一个只有 \(5\) 个球的小袋中，放回和不放回的差距大到不能忽视。

\textbf{第 4 步：独立还是依赖。}不同事件之间是否存在信息通道？扔两次硬币通常假设独立，但如果是同一阵风吹动了两次抛掷的手臂，它们可能就不独立。你不需要每次都证明独立，但你要意识到自己在什么时候调用了独立假设。独立是一个精确的数学条件，不是``看起来没关系''的等价词。

\textbf{第 5 步：观察机制。}你是直接看到了某个结果，还是通过一个可能出错的检测手段得知的？Monty Hall 问题里，主持人的开门规则决定了已知信息的内容和强度；医院检测问题里，假阳性率决定了``阳性''这个证据能带来多少信息。同样的已知事件，不同的产生机制会对应不同的似然，这在后面 Bayes 公式中会变得致命重要。

\subsection{什么时候必须小心隐变量}

你口袋里有两枚硬币：一枚公平（正面概率 \(0.5\)），一枚双正面（正面概率 \(1\)）。你闭眼摸出一枚，掷一次，正面。这枚硬币是双正面币的概率是多少？

你可能会想：反正已经掷出正面了，公平币也能掷出正面，是不是应该用 Bayes 公式？对——但这里有一个更微妙的陷阱：你拿出来的到底是一枚硬币还是另一枚，这个过程叫``随机选取''，而你一旦选定了，掷出的结果就是这枚硬币的固有属性决定的。样本空间里的结果，是``选了哪枚硬币''和``掷出哪一面''的配对。第一层随机（选币）和第二层随机（掷硬币）在物理上泾渭分明，在概率空间里却是同一个联合分布。

如果你错误地把样本空间只定义成\{正面, 背面\}，你会漏掉``选到哪枚硬币''这个隐变量。而一旦你意识到必须同时记录选币结果和掷币结果，样本空间自然变成
\[
\{(\text{公平, 正}),\ (\text{公平, 反}),\ (\text{双正面, 正})\},
\]
概率分别是 \(0.5\times0.5=0.25\)，\(0.5\times0.5=0.25\)，\(0.5\times1=0.5\)。在这个正确的样本空间上，已知正面后双正面币的概率是 \(0.5/(0.25+0.5)=2/3\)。

这个例子小，但暴露的建模问题很大：现实世界的``一次实验''可能包含多个子过程（选币 + 掷币），你的样本空间必须把每个子过程的结果都编码进去。漏掉一个子过程，就是漏掉了信息的来源，后面的条件化就失去了根基。

\subsection{模型可以有用而不必完美真实}

公平硬币、独立掷骰、Poisson 到达——现实中没有哪个严格满足这些假设。硬币有微小不均匀，掷骰的手臂动作不是完全不可预测的，顾客到达率在午餐时间和下午三点不可能一样。

但这不等于模型没用。一个模型好不好，不该只问``是不是严格真实''，而应该问：

\begin{itemize}
\tightlist
\item 它保留了原问题中哪些最关键的结构？（例如，Poisson 过程保留了``事件独立发生''和``速率稳定''这两个核心特征，忽略了顾客之间的社交影响。）
\item 被忽略的因素在什么条件下会严重改变结论？（高峰期的到达不再是稳定速率，排队估计需要换模型。）
\item 模型参数能不能通过数据校准？（如果校准后发现速率确实随时间大幅变化，你就有理由拒绝稳定速率的假设。）
\item 模型在什么范围内给出的是可靠近似，什么范围内可能给出方向性错误？
\end{itemize}

一个经典的例子：把大量独立顾客的到达近似为 Poisson 过程。在平稳时段（比如工作日下午两点到三点），这个近似往往给出不错的排队时间估计。但如果你用同一个模型去估计午餐高峰或者大型活动散场时的排队时间，结论可能错得离谱——不是因为 Poisson 过程的数学有问题，而是``稳定速率''这个建模假设在那个时间段不成立。

再举一个更日常的例子。你要估计明天降水的概率，天气预报说 \(30\%\)。这个 \(30\%\) 是什么意思？它不是一个硬币的固有概率——明天只有一次，要么下雨要么不下。它代表的是：在气象模型产生的类似初始条件的历史集合中，约 \(30\%\) 的日子下了雨。这个数字依赖模型的分辨率、历史数据的代表性、以及``类似''的定义。如果你把它当成一个客观概率像掷硬币一样使用，你实际上接受了一整套模型背后的假设。接受假设不是错——但不知道自己在接受假设，以后翻了车都找不到哪里翻了。

第一层严谨性：给定模型之后，数学推导必须没有错误。第二层严谨性：从现实场景到模型的翻译过程必须透明，假设必须写清楚。只完成第一层而回避第二层，会出现``计算过程挑不出毛病，但结论在现实中毫无意义''的尴尬。

\subsection{下一步：新信息会改写你的概率空间}

到目前为止，\(P(A)\) 描述的是当前模型下事件 \(A\) 的概率——一个固定不变的数。但你每时每刻都在接收新信息：检测报告出来了，第一张牌翻开了，前十分钟没有顾客到达。接收到新信息之后，旧的概率分配还能直接用吗？

先回到第一节那个两个孩子的例子。已知一个家庭有两个孩子，\(P(\text{两个男孩})\) 在没有任何附加信息时是 \(1/4\)（假设生男生女各半且独立）。但如果你知道了``至少有一个男孩''，原来的样本空间 \(\{BB,BG,GB,GG\}\) 缩小为 \(\{BB,BG,GB\}\)，概率变成了 \(1/3\)。更进一步，如果你知道``较大的孩子是男孩''，样本空间缩小为 \(\{BB,BG\}\)，概率变成 \(1/2\)。

你做的事是：把与观察不相容的那些结果删掉，然后在剩下的结果上重新做规范化，让它们的概率加起来回到 \(1\)。对原概率 \(P(A)\) 做加减乘除没法完成这个更新——更新需要一种新的运算：条件概率。

你可能会想：信息的强度可以量化吗？有些信息几乎不改变概率——比如你被告知路上有一个水坑，而你本来就打算穿雨鞋，这条信息对你的穿衣决策几乎无影响。有些信息戏剧性地翻转概率——体检报告上出现一个红色箭头。信息量的大小和信息改变概率的幅度之间有什么精确的关系？这恰好是信息论中熵与互信息的出发点，在\hyperref[ch:041]{``信息论：熵、编码与几何''}中会回到这条线索上来。

而眼下最紧迫的问题是：怎么从``原因到结果的概率''反向算出``结果到原因的概率''——因为你日常碰到的推断几乎都是反方向的。

这就是下一章\hyperref[ch:038]{``条件概率、Bayes 与独立性：新信息怎样改变判断''}的起点。你会看到，整个统计推断的骨架——从医学诊断到机器学习分类——都挂在条件概率和一张叫做 Bayes 公式的网上。
